\section{Background and Related Work}
\label{sec:related}

\subsection{From place embeddings to per-visit embeddings}
\label{sec:related-embeddings}
A place embedding transforms a POI into a vector so that similar or nearby
places sit close together in the vector space.
Deep Graph Infomax (DGI)~\cite{velickovic2019dgi}, a general self-supervised
method for graphs, learns such vectors on a place network linked by
similarity, time, and distance.
Hierarchical Graph Infomax (HGI) \cite{huang2023hgi} adds geographic structure,
organizing the network as place, then region, then city. Both produce one vector per place, so two visits to the same coffee shop look
identical to the model.
Contextual check-in representations instead give each visit its own vector: CTLE
\cite{lin2021ctle}, the closest example, learns a vector per visit by masking and
reconstructing parts of a user's check-in sequence. CTLE is a sequence model, a Transformer that reads the
check-in sequence itself; our representation, Check2HGI, remains a network model, keeping the
hierarchical place-to-region-to-city network and the same infomax objective,
now extended one level deeper, from individual places down to individual
check-ins (Fig.~\ref{fig:dataflow}). The order of a user's visits still enters
the network, as links between consecutive check-ins
(Section~\ref{sec:method-rep}).

We compare against CTLE and a feature-concatenation control to separate the
combination from contextualization alone and from feature injection.

\subsection{Predicting the next category and the next region}
\label{sec:related-tasks}
A large body of work predicts the next place that a user will visit, the exact POI,
with recurrent and attention models such as DeepMove \cite{feng2018deepmove}, STAN
\cite{luo2021stan}, and GETNext \cite{yang2022getnext}. In contrast, we target two properties of the next visit, its category and its
region, rather than the exact place. We choose them because they are
what a mobility-aware service can act on, not to make the task easier; the region
alone spans hundreds to several thousand classes. Predicting
over a partition of the map is standard in the human-mobility literature, with a grid cell as the
target~\cite{luca2021mobilitysurvey}; we substitute official neighborhood-scale units. Our earlier
work~\cite{silva2025mtlnet} paired static category classification with next-category prediction. This paper introduces the next-region task, adds the check-in-level
representation, and reports, for each task, where the joint model outperforms a dedicated single-task model
and where their difference stays within a stated bound (Section~\ref{sec:results-part2}).

The field
increasingly models several granularities at once; in those systems, category and
region are auxiliary signals that help a primary next-place task
(MCMG \cite{sun2024mcmg}, HMT-GRN \cite{lim2022hmtgrn}).
We study the pair as the object itself. To our knowledge, fine-grained region as an
end target of equal standing, rather than an auxiliary coarse grid cell, is
underexplored.
The nearest exceptions do not study our exact pairing:
DRRGNN~\cite{zhu2022drrgnn} forecasts a person's next activity region jointly
with a mobility-intention label, but over regions discovered per person rather
than a fixed citywide partition; a generative recommender~\cite{sun2025kgtb}
predicts category and region together, but only as auxiliary steps toward its
next-place ranking.

\subsection{Multi-task learning for mobility}
\label{sec:related-mtl}
MCARNN \cite{Liao2018} jointly predicts activity and
location, and a recent hierarchy-aware model continues the
pattern \cite{wang2025hamtl}; in both, the location target is
the exact place. The closest alternative to our setting is
the cascade: predict the category first, then the place given the
category \cite{ye2013nextmove}. CSLSL
\cite{huang2024cslsl} is its current form, predicting in a chain (when, then what,
then where), with location as the primary output and category as an
intermediate step; CatDM \cite{yu2020catdm} likewise uses the predicted
category only to filter candidate places. We instead predict
category and region in parallel, as end targets of equal standing in one model
with a single forward pass, and we drop the next-place target entirely. On optimization, we are conservative by design: joint training with a fixed loss weighting is
standard practice~\cite{caruana1997multitask}, and gradient-balancing methods (PCGrad~\cite{yu2020pcgrad}, Nash-MTL~\cite{navon2022nashmtl})
rarely improve on a well-tuned fixed weighting with two tasks~\cite{xin2022domtl}. We confirmed
this during development, on an earlier preparation of the data: a screen of nineteen loss and
gradient balancers at their default configurations, at a single seed on Alabama and Florida, found
none that improved on fixed equal weighting on both tasks at both datasets.  The cosine similarity between the next-category and next-region updates on the
shared trunk, measured on the reported joint model at four of the six datasets (Istanbul, Alabama,
Arizona, and Florida), is equivalent to zero at every one of them, with per-dataset means within
two thousandths of zero, so a balancer has no conflict to resolve. This is a finding for this pair
of tasks, not a general rule.

